\documentclass[9pt,twocolumn,twoside]{pnas-new}
% Use the lineno option to display guide line numbers if required.

\templatetype{pnasresearcharticle}
\setboolean{displaywatermark}{false}

\usepackage{siunitx}
\input{numbers}
\graphicspath{{../figures/}}

% Footer: course instead of journal placeholders (template class left untouched)
\makeatletter
\fancypagestyle{firststyle}{%
  \fancyfoot[R]{\footerfont 02805 Social Graphs and Interactions\hspace{7pt}|\hspace{7pt}\textbf{November 2026}\hspace{7pt}|\hspace{7pt}\textbf{\thepage\textendash\pageref{LastPage}}}%
  \fancyfoot[L]{\footerfont\@ifundefined{@doi}{}{\@doi}}}
\fancyfoot[RO]{\footerfont 02805 Social Graphs and Interactions\hspace{7pt}|\hspace{7pt}November 2026\hspace{7pt}|\hspace{7pt}\textbf{\thepage}}
\makeatother

\title{Many faces, few hands: how credit networks reveal Spotify's ghost artists, and where the recommender puts them}

\author[a]{August Holm J{\o}rgensen}
\author[a]{Bertram Nyvold Larsen}
\author[a]{Saxe Br{\o}ten}

\affil[a]{DTU Compute, Technical University of Denmark, 02805 Social Graphs and Interactions. Student IDs s224166, s224194, s224224.}

\leadauthor{J{\o}rgensen}

\significancestatement{Playlists decide much of what people hear while they work, study or sleep. On Spotify's mood playlists, most of the artist names are not independent musicians. They are labels on catalogues made by a handful of production companies. The credits hide this, but the codes that register each recording reveal it, and three independent measurements point to the same names.}

% ------------------------------------------------------------------ TODO before hand-in
% Replace with what each of us actually did. The course does not accept "all contributed equally".
\authorcontributions{Author contributions: TODO -- A.H.J., B.N.L. and S.B. to write one specific sentence each (e.g. who designed the three dimensions, who built the network nulls, who did the text analysis, who wrote which sections).}
\authordeclaration{The authors declare no competing interests. Code and processed data: \url{https://github.com/Benyla/02805_SGI} (folder \texttt{project/}).}
\correspondingauthor{Course project, hand-in 27 November 2026.}

\keywords{streaming $|$ ghost artists $|$ bipartite networks $|$ null models $|$ recommender audit}

\begin{abstract}
Spotify's mood playlists are said to be filled with cheap stock music released under invented artist names. We test this on one snapshot of \nPlaylists{} Spotify editorial playlists (Denmark, November 2026): \nTracks{} tracks by \nNames{} artist names, with their credits, recording codes (ISRC), \nFal{} ``Fans also like'' links, bios and cover art. We score every name on three independent dimensions: identity, production and demand. They agree (Spearman $\rho=\rhoPDF$--$\rhoIPF$ among mood-playlist artists). Of \npNames{} ghost-like names found \emph{without} production data, \npTopEight{} trace back to just eight registrants of recordings, seven of them Swedish. Ghost-like names hold \slotGhostF{} of mood-playlist slots. The recommender keeps them apart from real artists (\flowRGDeg$\times$ the expected links) and groups names from the same house (\samePl$\times$ chance, even within playlists). Names, titles and sound do not give a ghost away; bios do.
\end{abstract}

\dates{This manuscript was compiled on \today}
\doi{\url{https://github.com/Benyla/02805_SGI}}

\begin{document}

\maketitle
\thispagestyle{firststyle}
\ifthenelse{\boolean{shortarticle}}{\ifthenelse{\boolean{singlecolumn}}{\abscontentformatted}{\abscontent}}{}

% =====================================================================================
\dropcap{P}laylists with names like \emph{Deep Focus}, \emph{Sleep} or \emph{Jazz in the Background} are among the most followed on Spotify, and placement on them is worth millions of streams~\cite{aguiar2021platforms,eriksson2019teardown}. Journalists have reported that much of this music is not made by the artists it is credited to. Pelly describes an internal Spotify programme, ``Perfect Fit Content'', in which production companies supply cheap background music for editorial playlists~\cite{pelly2025harpers,pelly2025mood}. In 2022 the Swedish newspaper \emph{Dagens Nyheter} (DN) obtained a list of 830 artist names linked to one company, Firefly Entertainment; at least 495 of them were on Spotify's own playlists, and music for more than 500 of the names came from about 20 songwriters~\cite{ingham2022mbw}. A later DN story traced one composer behind more than 650 names, released through Overtone Studios, a label owned by Epidemic Sound~\cite{stassen2024mbw}. The question has become more pressing as generative AI lowers the cost of such catalogues: in July 2026 Deezer reported that more than half of its daily uploads were fully AI-generated~\cite{techcrunch2026deezer}.

This reporting rests on leaked lists and interviews. We ask whether the pattern can be measured from public data, using the tools of network science and text analysis. Our question is: \emph{do the many anonymous names on Spotify's mood playlists come from a few production houses, and where does Spotify's recommender put them relative to real artists?}

A ``ghost artist'' is not a single property, so we measure three. \emph{Identity}: is the name a recognisable person or band (verified, a bio, links, photos)? \emph{Production}: is the catalogue made by a company that also stands behind many other names? \emph{Demand}: do the streams come from fans, or from placement on a playlist? Each is measured with different data. If all three point to the same names, that agreement is our main check that we are measuring something real, and not an artefact of one data source. We then use the names flagged this way to answer the two network questions: how concentrated the ownership is (a bipartite name--company network), and how the ``Fans also like'' recommendations connect ghost-like names and real artists, compared with degree-preserving null models.

% =====================================================================================
\section*{Results}

\subsection*{The data}
Spotify's official interface no longer gives new apps access to editorial playlists or related artists~\cite{spotify2024api}, so we collected a single snapshot through the web player on 4 November 2026 (Materials and Methods). Figure~\ref{fig:data} shows how the pieces fit together. We chose \nPlaylists{} hand-curated Spotify playlists: \nFunctional{} \emph{functional} ones (sleep, focus, chill, mood jazz and classical) and \nGenre{} \emph{genre} playlists with known real artists (e.g.\ \emph{Rap UK}, \emph{ALT NOW}) as a control. Their \nSlots{} slots hold \nTracks{} tracks by \nNames{} artist names. For every track we have the credits and the ``Source'' (the label shown in the credits); for every name, its profile and its \nFal{} ``Fans also like'' links. We matched tracks to Deezer to get each recording's ISRC, an international code whose first five characters identify the country and the organisation that registered the recording. We call that organisation the \emph{house}. Deezer matched \deezerMatch{} of the tracks.

The contrast between the two kinds of playlist is stark. Artists on functional playlists have a median of \fplF{} followers per monthly listener, against \fplG{} on genre playlists: real fans follow the artists they listen to, while playlist listeners do not. \aliasOnlyF{} of functional tracks credit nobody but the artist name itself (\aliasOnlyG{} on genre playlists), and only \verifiedF{} of their artists are verified (\verifiedG).

\begin{figure*}[t]
\centering
\includegraphics[width=\linewidth]{fig1_data.pdf}
\caption{\textbf{The data as three linked layers.} Playlists (top) contain artist names (middle); each name's tracks are registered by an ISRC house (bottom). Three networks connect the layers: \emph{placement} (name is on a playlist), \emph{Fans also like} (Spotify recommends one name on another's page) and \emph{ownership} (a house registered the name's tracks). The diagram is a schematic; the counts on the right are the real ones. Text (titles, bios, playlist descriptions) and cover images are attached to the names.}
\label{fig:data}
\end{figure*}

\subsection*{Three independent views of a ghost agree}
We gave every name three scores, oriented so that higher means more ghost-like (Materials and Methods). The \emph{identity} score is the share of five public-identity markers that are missing. The \emph{production} score combines the size of the name's ISRC house with the share of its tracks that credit no real person. The \emph{demand} score combines few followers per listener, depending on one of our playlists as the main place listeners find the name, and listeners spread thinly around the world. None of the scores uses the functional/genre label.

The three scores rise together (Fig.~\ref{fig:convergence}A). To make sure this is not just the difference between the two kinds of playlist, we computed the correlations among functional-playlist artists only: identity and production $\rho=\rhoIPF$ \rhoIPFCI, identity and demand $\rho=\rhoIDF$ \rhoIDFCI, production and demand $\rho=\rhoPDF$ \rhoPDFCI\ (95\% bootstrap intervals). Among the genre-playlist artists, who are almost all real, the scores do not correlate ($\rho$ between \rhoPDG{} and \rhoIPG), which is what we expect if the scores pick up ghosts rather than noise.

We turned the three scores into one label with a two-component Gaussian mixture model. \nGhost{} of the \nNames{} names are \emph{ghost-like}; only \uncertainShare{} fall in the uncertain middle ($0.1<p<0.9$). We checked the label on groups it never saw (Fig.~\ref{fig:convergence}B). \refFirefly{} of the \nFirefly{} names whose tracks carry Firefly Entertainment, the company in the DN investigation, are ghost-like. The \nMajors{} names from classical major labels (Deutsche Grammophon, Decca, Sony, Warner), which sit on the \emph{same} functional playlists, are at most \refMajorsMax{} ghost-like, and the genre-playlist artists \refGenre. One known supplier is missed: names released by Epidemic Sound itself are classified real (\refEpidemic, $n=\nEpidemic$). Their profiles have bios, links and photos, and \epidemicFplRatio$\times$ the followers per listener of the median functional artist. A supplier that invests in its artists' identities is invisible to the identity dimension.

\emph{What the label captures and misses.} It captures names that look, are produced and are listened to like stock catalogue. It is a measurement, not a verdict about any person: some ghost-like names are pseudonyms of real working musicians.

Ghost-like names hold \slotGhostF{} of all functional-playlist slots and none of the genre slots (Fig.~\ref{fig:convergence}C). Lifetime plays of their tracks on these playlists add up to \playsGhost{} billion, against \playsReal{} billion for the real artists beside them. Only a few functional playlists are built around real artists, such as \emph{Calming Classical} and \emph{Women of Ambient}.

\begin{figure*}[t]
\centering
\includegraphics[width=\linewidth]{fig2_convergence.pdf}
\caption{\textbf{Three views agree, and the label passes checks it never saw.} (\emph{A}) Functional-playlist artists ($n$=4,530) grouped by how many of five identity markers they lack; mean production and demand scores rise with it (bands: 95\% bootstrap CI). Within these artists the Spearman correlations are \rhoIPF{} (identity--production), \rhoIDF{} (identity--demand) and \rhoPDF{} (production--demand); among genre-playlist artists they are near zero. (\emph{B}) Share labelled ghost-like in reference groups not used to build the label: the company named in the DN investigation (orange), the classical majors on the same playlists, and the genre-playlist artists. Epidemic Sound's own names are missed (see text). (\emph{C}) Each dot is a playlist (area $\propto$ followers); the dashed line is the share over all \nSlots{} slots.}
\label{fig:convergence}
\end{figure*}

\subsection*{Many faces, few hands}
The credits hide who makes the music. \selfSourceGhost{} of ghost-like names give \emph{their own name} as the Source of their tracks, against \selfSourceReal{} of the real artists on the same playlists. In the credits they look self-released. The ISRC tells a different story: self-labelled names sit in houses with a median of \selfSourceHouse{} other names.

Figure~\ref{fig:ownership}A draws the ownership network. Each dark node is a house and each dot around it a name whose tracks it registered, so the size of a disc shows how many names a house holds. Because the production score already uses house size, we colour the names here with a label built from identity and demand only (\nGhostNoprod{} names; it agrees with the full label on \labelAgreement{} of names). Even with this label, the \npNames{} ghost-like names that have an ISRC come from only \npHouses{} houses, and \npTopEight{} of them come from just eight (\npTopTen{} from ten). \npSE{} are Swedish-registered, against \seReal{} of the real artists. A ghost-like name typically shares its house with \npMedGhost{} other names; a real artist with \npMedReal{} (Fig.~\ref{fig:ownership}B). The same eight houses fill \npTopEightSlots{} of the ghost-like playlist slots (Fig.~\ref{fig:ownership}C). For four of the eight largest houses no company appears anywhere in the credits: every name is its own label. Two of the houses belong to Overtone Studios, which is owned by Epidemic Sound~\cite{stassen2024mbw}; this is how the Epidemic catalogue shows up even though its own-brand names are classified real. \nMultiHouse{} names have tracks in more than one house (faint arcs in Fig.~\ref{fig:ownership}A), linking the large houses to each other.

\emph{What this captures and misses.} An ISRC registrant is the organisation that registered the recording. It can be a label or a distributor, not necessarily the composer. A house is therefore an upper bound on common control, and the real number of people behind the names may be smaller still~\cite{ingham2022mbw}.

\begin{figure*}[t]
\centering
\includegraphics[width=\linewidth]{fig3_ownership.pdf}
\caption{\textbf{Many faces on few hands.} (\emph{A}) Ownership network of the \nWithRegistrant{} names with an ISRC. Each dark node is a house (ISRC registrant); the names whose tracks it registered sit on a spiral around it, so disc area is proportional to the number of names. Orange: ghost-like by identity and demand alone (no production data). Labels give the registrant code, the number of names and the company in the credits, if any. (\emph{B}) Seen from a name, how large its house is (complementary cumulative distribution). (\emph{C}) Cumulative share of functional-playlist slots covered by the largest houses: eight houses cover \npTopEightSlots{} of the ghost-like slots, while real artists' slots are spread over hundreds of houses.}
\label{fig:ownership}
\end{figure*}

\begin{figure*}[t]
\centering
\includegraphics[width=\linewidth]{fig4_recommender.pdf}
\caption{\textbf{The recommender's map is a map of the playlists.} (\emph{A}) ``Fans also like'' links among the dataset's names (largest connected component, \nGiant{} names). Each Louvain community is laid out on its own; labels give its most common playlist, size and ghost-like share. (\emph{B}) Observed divided by expected number of links of each kind under two nulls (1,000 samples each). Open: degrees kept; filled: degrees and the target's home playlist kept. All deviations are significant ($|z|\geq\zRGPl$). (\emph{C}) NMI between the communities and three partitions of the names; hatched bars show shuffled labels.}
\label{fig:recommender}
\end{figure*}



\subsection*{Where the recommender puts them}
Every artist page has a list of 20 ``Fans also like'' artists. These lists are the recommender's view of who is similar to whom. \nFalIn{} of the links point to another name in our data. On real artists' pages, \realSlotsToGhost{} of all recommendation slots go to ghost-like names. This is a lower bound, because \inSetReal{} of real artists' links point outside our data.

To see whether ghosts are mixed in or kept apart, we counted the links between the two groups and compared them with a null model that keeps every name's number of incoming and outgoing links but shuffles who links to whom~\cite{maslov2002specificity} (Fig.~\ref{fig:recommender}B, open circles). Real artists link to ghost-like names \flowRGDeg$\times$ as often as expected, and ghost-like names to real artists \flowGRDeg$\times$ ($|z|=\zRGDeg$, 1,000 samples). We had expected a ``funnel'' in which real artists send listeners on to ghosts while ghosts rarely send them back. There is none: the two directions are depleted almost equally.

Much of this separation comes from the playlists. Listeners who found a name on \emph{Sleep} also listen to other \emph{Sleep} names. A second null keeps the degrees \emph{and} the home playlist of every link's target. Against it, the separation shrinks to \flowRGPl$\times$ and \flowGRPl$\times$ (filled circles; $|z|=\zRGPl$). Names from the same house, however, are still recommended next to each other \samePl$\times$ as often as expected (\sameObs{} links against \sameExpPl; $z=\zSamePl$). The recommender groups a house's names together even beyond what the playlists explain.

Community detection gives the same picture. The Louvain method~\cite{blondel2008louvain} splits the network into \nComm{} communities with modularity $Q=\louvainQ$, far above the $Q=\louvainQnull\pm\louvainQnullSd$ of degree-preserving rewired networks~\cite{guimera2004modularity}. The communities follow the playlists: their normalised mutual information (NMI)~\cite{danon2005nmi} with each name's home playlist is \nmiHome{} (shuffled: \nmiHomeNull), against \nmiHouse{} with the house (shuffled: \nmiHouseNull) (Fig.~\ref{fig:recommender}A, C).

\emph{What this captures and misses.} ``Fans also like'' is the output of the recommender, not its inner workings, and modularity measures how grouped a partition is, not whether it is right. The result says that the recommender's map of similarity is largely a map of Spotify's own playlists. Placement creates co-listening, and co-listening becomes ``similarity''.



\subsection*{What a listener sees}
\begin{figure}[!tb]
\centering
\includegraphics[width=\linewidth]{fig5_text.pdf}
\caption{\textbf{Ghosts are hard to see from the outside.} (\emph{A}) Cross-validated ROC AUC for predicting the ghost-like label from one kind of content (functional-playlist artists). Hatched: 95\% range of the same classifier on shuffled labels. *Exploratory. (\emph{B}) Words most typical of ghost-like and real bios (weighted log-odds with an informative prior~\cite{monroe2008fightin}), and how often a bio names an entity.}
\label{fig:text}
\end{figure}

Would a listener notice? We trained simple classifiers to recognise ghost-like names from one kind of content at a time, among functional-playlist artists only (Fig.~\ref{fig:text}A). The artist name (character n-grams) gives a cross-validated ROC AUC of \aucName{}, and track titles \aucTitles{}, where 0.5 is guessing. Ghost-like names look like people: \personGhost{} are a two-word ``Firstname Lastname'', just like the real artists beside them (\personRealF). As exploratory checks, cover art (CLIP embeddings~\cite{radford2021clip}) reaches \aucCover{}, and the sound of 30-second previews (\nAudio{} tracks) \aucAudio{}: on the same playlists, ghost tracks sound like real ones. Covers carry a house style, though: a cover's most similar cover under \emph{another} name comes from the same house \coverSameHouse{} of the time (shuffled: \coverSameHouseNull).

Bios are the exception (AUC \aucBio, $n=\nBio$), but only \bioGhost{} of ghost-like names have a bio at all, against \bioReal{} of real artists. Ghost bios describe a mood (\emph{calm}, \emph{gentle}, \emph{stillness}, \emph{designed}); real bios describe a career (\emph{album}, \emph{released}, \emph{orchestra}, \emph{debut}) (Fig.~\ref{fig:text}B)~\cite{monroe2008fightin}. Named-entity recognition finds a person, place, date or organisation in \entReal{} of real bios and \entGhost{} of ghost bios, which are also four times shorter (median \tokGhost{} against \tokReal{} words).



% =====================================================================================
\section*{Discussion}
Three independent views (who the name is, who makes the music, and where the listening comes from) point to the same names, and they pass checks they never saw. Using them, we find that two-thirds of the slots on Spotify's Danish mood playlists go to names whose recordings come from a small number of mostly Swedish houses. The credits hide this through self-labelling, and only the recording codes reveal it. The recommender does not funnel listeners from real artists to ghosts. Instead, it mirrors the playlists and groups each house's names together. A listener can hardly tell: names, titles and sound look the same, and only the rare bio gives a ghost away.

\emph{Limitations.} This is one snapshot from one market, so we cannot see how playlists change over time, for example whether ghosts displace real artists. Daily snapshots would allow that. The ghost-like label is an operational definition. It misses suppliers who build credible profiles (Epidemic Sound's own names), and its cut-off comes from a mixture model. Deezer matching is by artist and title, and it fails more often for functional tracks (\deezerMatchF{} against \deezerMatchG), so houses are measured on a subset. An ISRC registrant can be a distributor. Finally, ``Fans also like'' is one output of Spotify's recommender, and we observe it only once.
% TODO hand labels: when notes/hand_label_sheet.csv is filled in, run `python -m src.features.kappa`
% and add one sentence here with the inter-rater kappa and the model-vs-majority kappa.

\emph{Ethics.} Spotify's terms forbid automated collection, so we kept the crawl small and slow. We downloaded no audio from Spotify, and we used only Deezer's official previews, processed locally. Because some ghost names belong to real musicians, we report houses and companies in aggregate and never profile named individuals.

% =====================================================================================
\section*{Materials and Methods}
\small
\textbf{Collection.} On 4 November 2026 we queried Spotify's web-player API (\texttt{api-partner.spotify.com/pathfinder}) from an anonymous session in Denmark, at about 10 requests/s (22,027 requests). Of the 2,622 Spotify-owned playlists on the Danish browse page, we kept editorial playlists that are not filled per user by an algorithm (so every listener sees the same list) and chose \nFunctional{} functional and \nGenre{} genre playlists among them. For each track we stored credits and Source; for each album the (P) line and cover (640\,px); for each artist the profile, top cities, ``Fans also like'' (up to 20) and ``Discovered on'' lists. Tracks were matched to Deezer's public API by artist and title search with a duration check, giving the ISRC. The house is the first five ISRC characters (country and registrant code).

\textbf{Scores.} \emph{Identity} = share missing of: verified, bio, external link, gallery image, and a Spotify ``This Is'' playlist. \emph{Production} = mean of z-scored $\log_{10}$ house size (distinct primary-artist names registered by the name's dominant registrant) and the share of the name's tracks whose credits name only artists. The Source/(P)-line company size is left out because self-labelling hides it (see Results). \emph{Demand} = mean of z-scored $-\log_{10}$((followers+1)/(listeners+1)), whether the top ``Discovered on'' playlist is one of our \nPlaylists, and minus the listener share of the top city. Correlations are Spearman with 1,000 bootstrap resamples. The label is the posterior $p>0.5$ of a two-component Gaussian mixture on the three scores (or, for ownership, on identity and demand only), with a fixed seed.

\textbf{Networks.} \emph{Ownership}: bipartite graph of names and houses, with an edge when a house registered at least one of the name's tracks. \emph{Recommendation}: directed ``Fans also like'' edges between dataset names (\nFalIn). The degree null permutes the targets of all edges (a directed configuration model that keeps every in- and out-degree). The playlist null permutes targets only among edges whose target has the same home playlist (the playlist where the name has the most slots). Each null uses 1,000 samples, and we report $z=(\mathrm{obs}-\mu)/\sigma$. Communities come from Louvain on the undirected simple graph (igraph~\cite{csardi2006igraph}). The modularity null is 50 degree-preserving rewirings, and NMI is compared with 100 label shuffles. Layouts use Fruchterman--Reingold~\cite{fruchterman1991} inside each community or house; they are for display only.

\textbf{Text, images and sound.} Classifiers are TF-IDF (character 2--4-grams for names; word 1--2-grams for titles and bios) with logistic regression, using 5-fold stratified cross-validation and 20 shuffled-label runs as a null~\cite{pedregosa2011sklearn}. Bios had the artist's own name removed. Entities come from spaCy \texttt{en\_core\_web\_sm}~\cite{honnibal2020spacy}. Bio words are ranked by weighted log-odds with an informative Dirichlet prior~\cite{monroe2008fightin}. Cover embeddings come from OpenCLIP ViT-B/32 trained on LAION-2B~\cite{cherti2023openclip}, averaged per artist. Audio features (tempo, loudness, dynamic range, spectral centroid, chroma, MFCC summaries) were computed with librosa~\cite{mcfee2015librosa} on \nAudio{} Deezer previews: up to 8 ghost-like and 8 real tracks from each of \nAudioPlaylists{} playlists.

\textbf{Reproducibility.} Every number in this paper is written by \texttt{src/paper\_numbers.py} from the processed data, and every figure is made by one script in \texttt{src/figures/}. Random seeds are fixed.

\bibliography{refs}

\end{document}
